\documentclass[11pt]{article}
\usepackage[margin=1.05in]{geometry}
\usepackage{amsmath,amsthm,amssymb,mathtools,bm}
\usepackage{booktabs,array,enumitem}\usepackage[expansion=false]{microtype}
\usepackage{xcolor}
\usepackage[colorlinks=true,linkcolor=blue!60!black,citecolor=blue!60!black,urlcolor=blue!60!black]{hyperref}
\theoremstyle{plain}
\newtheorem{theorem}{Theorem}[section]\newtheorem{proposition}[theorem]{Proposition}\newtheorem{lemma}[theorem]{Lemma}
\newtheorem{corollary}[theorem]{Corollary}\newtheorem{conjecture}[theorem]{Conjecture}
\theoremstyle{definition}\newtheorem{definition}[theorem]{Definition}\newtheorem{remark}[theorem]{Remark}
\newtheorem{example}[theorem]{Example}
\newcommand{\R}{\mathbb R}\newcommand{\Z}{\mathbb Z}\newcommand{\E}{\mathbb E}\newcommand{\Var}{\operatorname{Var}}\newcommand{\Cov}{\operatorname{Cov}}
\newcommand{\relu}{\operatorname{relu}}\newcommand{\diag}{\operatorname{diag}}\newcommand{\tr}{\operatorname{tr}}
\newcommand{\Tr}{\operatorname{Tr}}\newcommand{\1}{\mathbf 1}\newcommand{\cA}{\mathcal A}\newcommand{\cB}{\mathcal B}
\newcommand{\cF}{\mathcal F}\newcommand{\cT}{\mathcal T}\newcommand{\cH}{\mathcal H}\newcommand{\aff}{\mathfrak{aff}}
\newcommand{\He}{\mathrm{He}}\newcommand{\Fa}{\Phi_{\mathrm F}}\newcommand{\cZ}{\mathcal Z}\newcommand{\Sym}{\operatorname{Sym}}
\newcommand{\Aff}{\mathrm{Aff}}\newcommand{\Ryshkov}{\mathcal R}
\newcommand{\Ap}{\mathsf A}\newcommand{\J}{\mathsf J}\newcommand{\cG}{\mathcal G}
\title{\textsc{Localization in the commutant, stage 15:}\\ \large deformed products, quasi-free lifts and the error cocycle.\\
The estimate as a renormalized character, the network as a random circuit,\\ and what a theory-native estimator must compute}
\author{}
\date{10 October 2026}
\begin{document}
\maketitle
\begin{abstract}
This stage starts from the problem, not from the chain. It takes two papers at their core and asks what they say about computing
$\E_x f_W(x)$ for one sampled network $W$.
\begin{itemize}[leftmargin=*]
\item Ebrahimi-Fard, Patras, Tapia and Zambotti (arXiv:1710.00735): a law is a deformation of the product, Wick polynomials are the
action $\mu^{-1}\star\mathrm{id}$, and renormalization is centering by a character of the noise.
\item Zhou (arXiv:2203.15920), after Kuan: a classical interacting process with a reflecting wall is the commutative
(Gelfand--Tsetlin) shadow of a free non-commutative random walk $P_t=(\mathrm{id}\otimes\langle\cdot\rangle_t)\Delta$.
\end{itemize}
Both run on one mechanism, a dynamics as convolution with a family of functionals through a coproduct. We read the network with it,
and with the quantum-simulation literature, where the same questions have sharp answers.
\begin{itemize}[leftmargin=*]
\item \emph{Wall jets (proved; checked).} Under the \emph{true} law of a pre-activation, the Wick (Appell) coefficients of \textsc{ReLU}
are the jets of the density at the wall: the transverse measure of the leaf and its normal derivatives. The joint cumulants of
\textsc{ReLU} outputs are sums of diagrams with wall-jet vertices and joint-cumulant hyperedges, no hyperedge inside one vertex. On a
non-Gaussian pair this expansion is $30\times$ more accurate than the Gaussian closure. Renormalizing the jets without the hyperedges
is \emph{worse} than not renormalizing at all.
\item \emph{Rows are random gates (proved).} The rows of a Gaussian weight matrix diagonalize Hermite degree exactly (Mehler). This is
the classical counterpart of Pauli-path orthogonality under random gates, and the radial (dilation) mode is its only obstruction.
\item \emph{The error cocycle (proved to leading order; measured on official network~0).} Every estimator's error splits into a
coherent modular part and an incoherent sum over layers of row-births transported by the first jet $\Phi(r)$. The sum rule
$\mathrm{MSE}\approx\sum_l\mathrm{birth}_l\times\mathrm{gain}_l$ closes to $4\%$ for the Gaussian closure once its collective mode is
removed (that mode is $61\%$ of its error), and to $10\%$ for the exact first-order chain, where the collective mode is already gone.
$79$--$95\%$ of the final error is inherited mean error. $81\%$ of the exact chain's error is born in layers $11$--$16$.
\item \emph{No free lunch for local corrections (proved).} Given the computed state, a birth's anisotropic part has zero conditional
mean. Only its isotropic part is predictable from per-neuron features. This explains stage~9's failed feature regression and bounds
what counterterms can do.
\item \emph{The quasi-free lift.} The fold $|z|$ is the reflecting wall (type~B), and births come only from the fold. The gate gas is the
image of the quasi-free Gaussian field under the sign map: arcsine pairs (Grothendieck--Krivine rounding) plus star cumulants
$-(4/\pi^2)\sum_i\prod_{j\neq i}\rho_{ij}$ (checked to $1$--$7\%$), the same stars as births.
\item \emph{Complexity.} The Gaussian backbone is the Clifford (matchgate) part, folds are magic, and a fold injects $M_\infty=0.188$
units of non-Gaussian energy per unit of field, so $n\,M_\infty/\sqrt{2\pi(1+t)}$ per layer. Generic quantum mean estimation would need
about $5\times$ the closures' budget for all outputs, so the resource being exploited is structure.
\end{itemize}
From these we derive what a theory-native estimator computes: an exact modular sector, an exact linear cocycle, and births computed
where birth $\times$ gain is large, decided per instance. Checks: \texttt{scripts/verify\_stage15.py}.
\end{abstract}
\tableofcontents
\section{The picture in one page, and what this stage refuses to assume}
\label{sec:summary}
\subsection{What this stage refuses to assume}
Each earlier stage ended with ``what this changes for the chain''. That framing assumes the answer: a system that propagates a truncated
law from layer to layer, plus corrections suggested by the theory. This stage drops four assumptions.
\begin{enumerate}[leftmargin=*]
\item \emph{That an estimator is a law propagator.} We first ask what \emph{any} estimator's error is made of (Section~\ref{sec:rows}).
\item \emph{That corrections are fitted.} We prove which parts of an error can be predicted from cheap features and which cannot
(Corollary~\ref{cor:nofreelunch}).
\item \emph{That the weights enter only through matrix products.} They are the noise of a regularity structure. Their ensemble expectations
are counterterms, and the information about the instance lives in the centred remainder (Section~\ref{sec:bphz}).
\item \emph{That the system is the same for every instance.} The error cocycle tells, per network, layer and row, where accuracy is born
and where it is lost (Section~\ref{sec:system}).
\end{enumerate}

\subsection{One mechanism in two papers, and in the network}
\begin{center}\small
\begin{tabular}{>{\raggedright\arraybackslash}p{3.6cm}>{\raggedright\arraybackslash}p{3.3cm}>{\raggedright\arraybackslash}p{3.0cm}>{\raggedright\arraybackslash}p{4.6cm}}\toprule
EFPTZ (1710.00735) & Zhou, Kuan (2203.15920) & quantum simulation & network\\\midrule
$\phi_\lambda=\lambda\star\mathrm{id}=(\lambda\otimes\mathrm{id})\Delta$ & $P_t=(\mathrm{id}\otimes\langle\cdot\rangle_t)\Delta$ & a channel & heat / OU semigroup on inputs; the layer map on laws\\
moment functional $\mu$ & the state $\langle\cdot\rangle_t$ & a state & the law of $z_l$\\
cumulants $\log^\star\mu$ & $t^{|\pi|}\prod\Tr$ (Faà di Bruno) & connected correlators & joint cumulants: the hyperedges\\
Wick map $\mu^{-1}\star\mathrm{id}$ & --- & normal ordering & Appell polynomials of the true law; the ReLU's coefficients are wall jets\\
deformed product $\cdot_\mu$ & --- & interaction picture & products of neurons under their true law\\
BPHZ character $\hat\mu^{-1}$ & --- & average over random gates & ensemble expectation over $W$: the annealed counterterm\\
centred model & --- & --- & the quenched, instance-specific part\\
restriction to a commutative subalgebra & Gelfand--Tsetlin subalgebra & computational basis & the gate (cell) subalgebra of the Gaussian field\\
reflecting wall (type B/D) & --- & --- & the fold $|z|$, Weyl-chamber quotient\\
--- & --- & Pauli weight & jet order (Hermite degree at the wall)\\
--- & --- & random-gate orthogonality & Mehler diagonalization by Gaussian rows\\
--- & --- & noise damps high weight & higher jets live only near walls\\
--- & --- & Clifford/Gaussian part & the Gaussian backbone $(\mu_l,C_l)$\\
--- & --- & magic & folds, $M_\infty=0.188$ per unit field\\\bottomrule
\end{tabular}
\end{center}

\subsection{The statements}
\begin{enumerate}[leftmargin=*]
\item \textbf{Wall-jet theorem} (Theorem~\ref{thm:walljet}; proved, checked). Under the true law, the Appell coefficients of
$\relu(z-c)$ are $\J_1=P(Z>c)$ and $\J_k=(-1)^kp^{(k-2)}(c)$: the density of the transverse measure at the wall and its normal
derivatives. Joint cumulants of \textsc{ReLU} outputs are diagram sums with jet vertices and cumulant hyperedges.
\item \textbf{Consistency} (Corollary~\ref{cor:consistency}; checked). The jets and the hyperedges belong to one deformation. Keeping
one without the other is worse than the Gaussian closure: on the test pair, $-2.7\times10^{-2}$ against $-1.2\times10^{-2}$, against
$+4\times10^{-4}$ with both.
\item \textbf{Rows are random gates} (Theorem~\ref{thm:rowmehler}; proved). This is the classical form of Pauli-path orthogonality. The
dilation is its only obstruction.
\item \textbf{The error cocycle} (Theorem~\ref{thm:cocycle}; leading order, measured). Error $=$ coherent modular part $+$ incoherent sum
of row-births transported by $\Phi(r)$. The measured sum rules close to $4\%$ and $10\%$.
\item \textbf{No free lunch} (Corollary~\ref{cor:nofreelunch}; proved). Only isotropic births are predictable from per-neuron
features.
\item \textbf{Quasi-free lift} (Section~\ref{sec:lift}). The fold is the reflecting wall; births come only from it; the gate gas is arcsine
pairs plus stars (checked to $1$--$7\%$).
\item \textbf{Complexity} (Section~\ref{sec:complexity}). The backbone is Clifford-like and the folds are magic, at density
$M_\infty/\sqrt{2\pi(1+t)}$ per neuron (checked to $1$--$9\%$). Quantum mean estimation is about $5\times$ worse than the closures for
all outputs.
\end{enumerate}

\subsection{The measured anatomy of an error (official network 0, truth from $10^9$ samples)}
\begin{center}\small
\begin{tabular}{>{\raggedright\arraybackslash}p{6.3cm}>{\centering\arraybackslash}p{3.4cm}>{\centering\arraybackslash}p{3.9cm}}\toprule
& Gaussian closure & exact first-order chain (+ radial $\kappa_4$)\\\midrule
final MSE & $4.06\times10^{-6}$ & $3.53\times10^{-8}$\\
share inherited from the previous layer's means (first jet) & $0.953$ & $0.790$\\
collective (dilation) share of the inherited error, layers $3$--$16$ & $0.48$--$0.68$ & $0.00$--$0.03$\\
cocycle sum rule $\sum_l\mathrm{birth}_l\,\mathrm{gain}_l$ vs MSE & $2.59$ vs $4.06\ (\times10^{-6})$ & $3.20$ vs $3.53\ (\times10^{-8})$\\
same, collective mode removed & $1.50$ vs $1.57\ (\times10^{-6})$ & $3.22$ vs $3.52\ (\times10^{-8})$\\
per-layer transport factor of mean error & $0.65$--$1.04$ & $0.65$--$0.94$\\
share of the MSE born in layers $11$--$16$ & $0.51$ & $0.81$\\\bottomrule
\end{tabular}
\end{center}
Three readings.
\begin{itemize}[leftmargin=*]
\item The organisers' radial $\kappa_4$ and the $\kappa_3$ memory remove exactly the coherent part.
\item What remains is a random walk of row-births across layers, dominated by the late layers.
\item No global correction can touch it; only better births can.
\end{itemize}

\section{Deformed products: the law as a character, the ReLU's Wick coefficients as wall jets}
\label{sec:hopf}
\subsection{The essence of Ebrahimi-Fard--Patras--Tapia--Zambotti}
Let $H=\R[x_a:a\in A]$ with the coproduct that makes every $x_a$ primitive, $\Delta x_a=x_a\otimes1+1\otimes x_a$. The dual product is
$(\alpha\star\beta)(P)=(\alpha\otimes\beta)\Delta P$, and the unital functionals $G(H)=\{\lambda:\lambda(1)=1\}$ form a group under $\star$. The paper's
points, in the form we use:
\begin{enumerate}[leftmargin=*]
\item \emph{Moments and cumulants are a group element and its logarithm}: $\mu=\exp^\star(\kappa)$ is the Leonov--Shiryaev relation, and the
cumulants form an infinitesimal character.
\item \emph{A law is a deformation of the product.} The Wick map $W=\mu^{-1}\star\mathrm{id}$ is the unique map with $W(1)=1$, $\partial_a W=W\partial_a$
and $\mu\circ W=\varepsilon$ on non-constants. It is a Hopf isomorphism onto $(H,\cdot_\mu,\Delta_\mu,\varepsilon_\mu=\mu)$, with
$:\!A_1A_2\!:\;=\;:\!A_1\!:\cdot_\mu:\!A_2\!:$. In the deformed algebra the expectation is the counit.
\item \emph{Renormalization is centring by a character of the noise.} On a comodule of trees over a Hopf algebra of forests, the BPHZ
character $\hat\mu^{-1}=\hat\mu\circ\hat S$ is the unique one making every non-empty tree centred (their Theorem 9.1). A good deformation
must preserve the remaining structure.
\end{enumerate}
The Gaussian input is one such group element. The heat semigroup is $\gamma_t=\exp^\star(t\,q)$ with $q(x_ax_b)=\delta_{ab}$, so
$e^{t\Delta/2}=\phi_{\gamma_t}$ and Hermite polynomials are $\gamma^{-1}\star\mathrm{id}$ of monomials. Stage~9's Eldan path and merge constants
live in this one-parameter subgroup.

\subsection{A layer acts on the character group}
\begin{proposition}[proved]\label{prop:layeraction}
Let $\mu_l\in G(H_l)$ be the moment functional of $z_l$.
\begin{enumerate}[leftmargin=*]
\item A linear map $z'=Wz$ induces $W^*:x'_i\mapsto\sum_jW_{ij}x_j$, which is a morphism of Hopf algebras. Hence $\mu'=\mu\circ W^*$ and
$\kappa'=\kappa\circ W^*$: cumulants transform tensorially and no new ones are created.
\item A coordinatewise map $g$ induces an algebra map that is not a coalgebra map. It creates cumulants, by the diagram law of
Theorem~\ref{thm:walljet}.
\item The Gaussian closure is the projection $\mu\mapsto\exp^\star(\kappa_1+\kappa_2)$ after every layer, a truncation in the Lie algebra. Its
neglected part is the \emph{vacuum} $\nu=\exp^\star(\kappa_{\ge3})$, with $\mu=\gamma\star\nu$ and $\E_\mu P=\E_\gamma[\phi_\nu P]$: the Edgeworth
operator is the action of $\nu$.
\end{enumerate}
\end{proposition}
\begin{proof}
(1) $W^*$ maps primitives to primitives and is multiplicative, and $\log^\star$ commutes with Hopf morphisms. (2) $g^*(x_a)=g(x_a)$ is not
primitive. (3) $\star$ is commutative, so $\exp^\star$ is a homomorphism from $(H^*,+)$ to $(G,\star)$.
\end{proof}

\subsection{The ReLU under its true law: wall jets}
\begin{theorem}[wall jets and the diagram law; proved, checked]\label{thm:walljet}
Let $Z$ have all moments and a density $p$ that is smooth near $c$, and let $\Ap_k=W(Z^k)$ be its Appell (Wick) polynomials. Then
$\Ap_k'=k\Ap_{k-1}$, $\E\Ap_k=0$ for $k\ge1$, and every polynomial $f$ satisfies $f=\sum_k\E[f^{(k)}(Z)]\,\Ap_k/k!$ exactly. For
$f=\relu(\cdot-c)$ the coefficients are
\[
\J_0=\E\relu(Z-c),\qquad \J_1=P(Z>c),\qquad \J_k=\E\,\delta^{(k-2)}(Z-c)=(-1)^kp^{(k-2)}(c)\quad(k\ge2).
\]
These are the density of the pre-activation at its wall and its normal derivatives. In the language of stages~11 and~14, they are the
jet of the transverse measure of the leaf $\{z=c\}$. For a pair, with $\J^X$, $\J^Y$ taken at the respective walls,
\begin{gather*}
\Cov\big(f(X),g(Y)\big)=\sum_{p,q\ge1}\frac{\E f^{(p)}(X)\,\E g^{(q)}(Y)}{p!\,q!}\,\E[\Ap_p(X)\Ap_q(Y)],\\
\E[\Ap_p(X)\Ap_q(Y)]=p!\,q!\,[s^pt^q]\exp\Big(\sum_{a,b\ge1}\frac{\kappa_{ab}s^at^b}{a!\,b!}\Big).
\end{gather*}
Equivalently, $\E[\Ap_p\Ap_q]$ is the sum over partitions of the $p+q$ legs in which every block meets both vertices, with weight
$\prod\kappa(\text{block})$. For $m$ neurons, the joint cumulants of $f_i(X_i)$ are the connected diagrams with jet vertices and joint-cumulant
hyperedges, where no hyperedge lies inside one vertex.
\end{theorem}
\begin{proof}
The Appell properties are EFPTZ Theorem~1.1. The expansion of $f$ follows from uniqueness in the Appell basis, by matching
$\E f^{(k)}$. For $\relu$, $\relu''=\delta$. For the pair, $\sum_p s^p\Ap_p(X)/p!=e^{sX}/\E e^{sX}$, so $\E[:\!e^{sX}\!:\,:\!e^{tY}\!:]=\exp(K(s,t)-K(s,0)-K(0,t))$,
where $K$ is the joint cumulant generating function. Only mixed cumulants survive. Faà di Bruno gives the partitions.
\end{proof}
For a Gaussian pair, only $\kappa_{11}$ survives and $\E[\Ap_p\Ap_q]=\delta_{pq}p!\kappa_{11}^p$ (Mehler). The jets are then
$s^{1-k}\He_{k-2}(-\alpha)\varphi(\alpha)$, which is the Gaussian closure's arc-cosine series. Every \emph{chain} is a partial
resummation of Theorem~\ref{thm:walljet} around Gaussian jets.

\emph{Check} (Section~\ref{sec:verify}). Take $X=g_1+0.15(g_1^2-1)$ and $Y=g_2-0.1(g_2^2-1)$ with $\rho(g_1,g_2)=0.6$; then $\kappa_{30}=0.93$,
$\kappa_{22}=-0.10$. The diagram law for $\E[\Ap_p\Ap_q]$, $p,q\le7$, holds to $2\times10^{-13}$, and $\Cov(X^3-2X,Y^4+Y)$ to ten digits. For
$\Cov(\relu(X-0.3),\relu(Y+0.2))=0.168248$:
\begin{center}\footnotesize\setlength{\tabcolsep}{3pt}
\begin{tabular}{lcccccc}\toprule
order $\le$ & 1 & 2 & 3 & 4 & 5 & 7\\\midrule
jets + all hyperedges & $-4.7\times10^{-2}$ & $-8.9\times10^{-3}$ & $+1.6\times10^{-3}$ & $+8.2\times10^{-4}$ & $+3.7\times10^{-4}$ & $+4.0\times10^{-4}$\\
jets + $\kappa_{11}$ only & $-4.7\times10^{-2}$ & $-2.6\times10^{-2}$ & $-2.7\times10^{-2}$ & $-2.7\times10^{-2}$ & $-2.7\times10^{-2}$ & $-2.7\times10^{-2}$\\\bottomrule
\end{tabular}
\end{center}
The Gaussian closure errs by $-1.2\times10^{-2}$. The full expansion is asymptotic. Its floor of about $4\times10^{-4}$ is set by the nearest
singularity of the marginal density, the fold caustic of $g\mapsto g+0.15(g^2-1)$, at distance $2.1$ from the wall.

\begin{corollary}[consistency of a deformation; checked]\label{cor:consistency}
The jets and the hyperedges are two faces of one deformation, the Wick map of the true law. Renormalizing one without the other changes
the product without changing the counit, and there is no reason for it to help. On the test pair it doubles the Gaussian closure's
error ($-2.7$ against $-1.2\times10^{-2}$), while the consistent expansion reduces it $30$-fold. EFPTZ's demand that a deformation preserve
the remaining structure is not decoration.
\end{corollary}

\subsection{Weights as the noise: BPHZ for an estimator}\label{sec:bphz}
In regularity structures the model $\Pi_\varepsilon$ is built from a noise, the counterterm is the noise expectation, and the renormalized
model is the centred one. Our noise is the weights.
\begin{itemize}[leftmargin=*]
\item Let $F(W_l;U)$ be any quantity an estimator forms at layer $l$, polynomial in the fresh matrix $W_l$ given the upstream state $U$.
\item Its character value $\lambda_U(F)=\E_{W_l}[F\mid U]$ is the annealed counterterm.
\item $F-\lambda_U(F)$, with nested subtractions given by the twisted antipode, is the centred (quenched) model.
\end{itemize}
Stages~9 and~13 found that quenched classes win and annealed controls carry nothing. This is the BPHZ split: the annealed part is already
inside the Gaussian backbone as its mean field, and the instance's information lives only in the centred parts. Theorem~\ref{thm:rowmehler}
below says these centred parts are orthogonal across distinct row multi-indices. That is why their contributions add in quadrature.

\section{Rows are random gates: the error cocycle}
\label{sec:rows}
\subsection{Gaussian rows diagonalize Hermite degree}
\begin{theorem}[rows are random gates; proved]\label{thm:rowmehler}
Let $W$ have i.i.d.\ $N(0,\sigma^2)$ entries, independent of an upstream $\sigma$-algebra $U$. Let $a,b$ be $U$-measurable vectors with
$\rho=\cos\angle(a,b)$, and write $\hat a=a/(\sigma|a|)$, $\hat b=b/(\sigma|b|)$.
\begin{itemize}[leftmargin=*]
\item For one row $w$: $\E[\He_p(w\cdot\hat a)\He_q(w\cdot\hat b)\mid U]=\delta_{pq}\,p!\,\rho^p$.
\item Distinct rows are independent. Hence, for multi-indices $\alpha\ne\beta$ over the output neurons,
$\E[\He_\alpha(W\hat a)\He_\beta(W\hat b)\mid U]=0$.
\end{itemize}
If a fixed scale $c$ is used in place of $|a|$, as any basis chosen in advance must do, degree $2$ leaks into degree $0$ by
$\E[\He_2(w\cdot a/\sigma c)]=|a|^2/c^2-1$, and every even degree leaks similarly. On the sphere, that is with the homogeneity used, the
orthogonality is exact.
\end{theorem}
\begin{proof}
$(w\cdot a,w\cdot b)$ is centred bivariate normal with the stated correlation. Mehler's formula applies.
\end{proof}
This is the classical counterpart of the orthogonality of Pauli paths under random single-qubit gates. That orthogonality is what makes
truncated Pauli-path and Pauli-propagation estimates of random circuits accurate on average: their mean squared error is the weight of
the dropped paths. Here a Gaussian row plays the random gate. The one thing that couples paths is the norm of the activation vector,
which is the dilation (stages~12--14). \emph{Check:} at $\rho=0.873$, all $p,q\le3$ agree within $|z|\le1.4$ ($4\times10^5$ rows); the
radial leak is $0.2346$ predicted against $0.2364$ measured.

\subsection{The error cocycle}
Fix any estimator with layer means $\hat m_l$ and fields $r_l$. Let $\delta_l=\hat m_l-m_l$ be its error in $\E h_l$. Define the first-jet
transport $T_l\delta=\Phi(r_l)\odot(W_l\delta)$ and the \emph{birth} $\beta_l:=\delta_l-T_l\delta_{l-1}$. Then, exactly,
\[
\delta_L=\sum_{l}T_{L}\cdots T_{l+1}\,\beta_l .
\]
\begin{theorem}[the cocycle; leading order in the law error]\label{thm:cocycle}
\begin{enumerate}[leftmargin=*]
\item \emph{Births live at walls.} By Theorem~\ref{thm:walljet} applied to $z=w\cdot h_{l-1}$,
\[
\beta_{l,i}=\sum_{k\ge2}\frac{\J_k(r_{l,i})}{k!}\,\big\langle\delta\kappa_k(h_{l-1}),\,w_{l,i}^{\otimes k}\big\rangle+\dots
\]
Here $\J_k=s^{1-k}\He_{k-2}(-r)\varphi(r)$ is concentrated near $r=0$. Only the first jet $\Phi(r)$, which transports, is spread out.
\item \emph{Births are i.i.d.\ over rows} given the upstream state, being functions of the fresh row $w_{l,i}$.
\item \emph{Births from different layers are orthogonal, except along the dilation.} Two transported births first meet at a fresh row.
Their cross term there is $\E[(w\cdot x)\beta(w)]=\sigma^2x\cdot\E\nabla\beta$ (Stein), which carries no Hermite-degree match and is
small. The exception is a dilation $\delta_{l-1}=\varepsilon m_{l-1}$. The first two jets together carry it to $\varepsilon(\Phi\mu+s\varphi)=\varepsilon m_l$
\emph{exactly} (homogeneity), so its births are coherent from layer to layer.
\item Hence $\mathrm{MSE}\approx\|\text{coherent modular part}\|^2+\sum_l\|T_{L}\cdots T_{l+1}\beta^\perp_l\|^2/n$.
\end{enumerate}
\end{theorem}
\begin{proof}[Sketch]
(1) Expand $\E\relu(w\cdot h)$ to first order in the difference between the estimated and true laws of $h_{l-1}$. The first jet is
removed into $T_l$. (2) Rows are i.i.d.\ given $U$. (3) Stein's lemma at the first common fresh layer; homogeneity of $\relu$ for the
dilation. (4) Expand the square and use (2)--(3).
\end{proof}
The theorem is the classical form of the random-circuit result. Dropping any piece of an estimator costs, additively, the energy of the
births it no longer prevents, times their downstream gain. Stage~9's single-age ladder and its finding that ``ages do not cancel''
(cumulative norm linear) are instances.

\subsection{Measured on official network 0}
Truth is from $10^9$ samples at every layer; its noise is about $8\times10^{-11}$ at the output.
\begin{center}\small
\begin{tabular}{lcccccccccccccccc}\toprule
layer & 1&2&3&4&5&6&7&8&9&10&11&12&13&14&15&16\\\midrule
\multicolumn{17}{l}{Gaussian closure: birth $\times$ downstream gain ($10^{-7}$)}\\
& .00&.28&.75&1.1&1.3&1.5&1.8&2.0&2.1&1.9&2.5&2.4&1.9&2.2&2.3&1.9\\
\multicolumn{17}{l}{exact first-order chain: birth $\times$ downstream gain ($10^{-9}$)}\\
& .02&.01&.05&.13&.26&.50&.69&1.2&1.4&1.8&2.5&3.5&3.1&3.6&6.0&7.4\\\bottomrule
\end{tabular}
\end{center}
\begin{itemize}[leftmargin=*]
\item \textbf{First jet.} The first jet carries $95\%$ (Gaussian closure) and $79\%$ (exact chain) of the final MSE.
\item \textbf{Transport factors.} Per layer they are $0.65$--$1.04$ and $0.65$--$0.94$: near-critical.
\item \textbf{The collective mode.} It holds $48$--$68\%$ of the Gaussian closure's inherited error at every layer, against $0$--$3\%$
for the exact chain.
\item \textbf{Sum rule.} It misses for the Gaussian closure ($2.59$ against $4.06\times10^{-6}$) exactly because of the collective
mode. With that mode projected out at every layer it closes to $4\%$ ($1.50$ against $1.57\times10^{-6}$). For the exact chain it
closes to $10\%$ ($3.20$ against $3.53\times10^{-8}$).
\item \textbf{Where errors are born.} The exact chain's births grow with depth: $81\%$ of its MSE is born in layers $11$--$16$, and
$42\%$ in the last two.
\end{itemize}

\begin{corollary}[no free lunch for local corrections; proved]\label{cor:nofreelunch}
Given the upstream state and the row's field $r_i$, $\E[\beta_{l,i}\mid r_i,U]=\beta^{\rm iso}(r_i)$. This is a function of the field
determined by the isotropic ($O(n)$-invariant) part of the law error. The anisotropic remainder has conditional mean zero. It is
unpredictable from the field, from $s_i$, or from any per-neuron statistic of the chain, unless the estimator already knows the
row-projected law error, in which case it should have corrected the law.
\end{corollary}
\begin{proof}
Condition on the projection of $w_i$ on the mean direction. The orthogonal part of $w_i$ is a fresh isotropic Gaussian, and it enters
only through $\langle\delta\kappa_k,w^{\otimes k}\rangle$, whose conditional mean is a trace.
\end{proof}
At the last layer, after removing the first jet, the part of the error explained by any smooth function of the field is $19\%$ (Gaussian
closure) and $7\%$ (exact chain). This settles three earlier findings:
\begin{itemize}[leftmargin=*]
\item stage~9's ridge regression on $34$ local features failed across networks;
\item v56's $103$ counterterms are worth $-3.8\%$ and no more;
\item the isotropic part is the modular sector, which should be \emph{computed} (stages~12--14), not fitted.
\end{itemize}

\subsection{Lattice reading}
Stage~14 predicted, by analogy with Rogers' second-moment formula, that the rows are independent given the state, correlated only through
the shared collective mode. The cocycle sum rule is the measured form of that statement. It holds once the collective mode is removed
and fails before.

\section{The quasi-free lift: reflecting walls, folds and the gate gas}
\label{sec:lift}
\subsection{The essence of Zhou's paper (after Kuan)}
Zhou studies a random surface growth in the quarter plane: interlacing particles that push and block, with a reflecting wall at $0$. It
is a one-parameter family of Plancherel measures for $O(\infty)$, so its characters are of type B/D. Three results matter here.
\begin{enumerate}[leftmargin=*]
\item \emph{The process is determinantal} along space-like paths, through a generalized $L$-ensemble.
\item \emph{It is the shadow of a free non-commutative walk.} On $U(\mathfrak{so}_{N+1})$ take the states $\langle X\rangle_\kappa=D(X)\kappa(I)$, with
$\kappa_t(O)=e^{t\Tr(O-I)}$. Their moments are an exponential of a trace,
\[\langle F_{i_1j_1}\cdots F_{i_mj_m}\rangle_t=\sum_\pi t^{|\pi|}\prod_B\Tr\prod_{b\in B}F_{i_bj_b},\]
which is Faà di Bruno again. Then
$P_t=(\mathrm{id}\otimes\langle\cdot\rangle_{\kappa_t})\circ\Delta$ is a semigroup that preserves the centre. Through the Harish-Chandra isomorphism, central
elements act by symmetric polynomials in the shifted labels, i.e.\ the particle positions. Restricted to the centre, $P_{t/2}$ is the
particle process of one level (Theorem~4.2). Restricted to the Gelfand--Tsetlin subalgebra, generated by the centres of the nested chain
$\mathfrak{so}_2\subset\dots\subset\mathfrak{so}_{N+1}$, it is the whole multi-level process (Theorem~4.3).
\item \emph{Gaussian free field, and the wall's type matters.} Moments of central elements converge to a three-dimensional Gaussian free
field. That field differs from the one of overlapping Wishart spectra, although both live to the right of a wall at $0$.
\end{enumerate}
The lessons are these.
\begin{itemize}[leftmargin=*]
\item[(L1)] A process with complicated classical interactions (blocking, pushing, reflection) is the commutative shadow of a free
non-commutative process. One computes upstairs.
\item[(L2)] The shadow is taken on a commutative subalgebra built from a nested chain.
\item[(L3)] What kind of wall it is (reflecting, or a hard edge) decides the fluctuations.
\end{itemize}
The mechanism is the one of Section~\ref{sec:hopf}: $P_t=(\mathrm{id}\otimes\langle\cdot\rangle_t)\Delta$ and $\phi_\lambda=(\lambda\otimes\mathrm{id})\Delta$ coincide for
cocommutative $\Delta$. A dynamics is convolution with a semigroup of functionals.

\subsection{The network's lift: the quasi-free field and its gate shadow}
For the network the upstairs object is the quasi-free (Gaussian) field of pre-activations, the backbone $(\mu_l,C_l)$.
\begin{itemize}[leftmargin=*]
\item The commutative subalgebra is generated by the gates, i.e.\ the cell indicators. This is the diagonal of stage~11's tile
groupoid.
\item These algebras increase with depth, matching $\mathrm{Gap}_l\subset\mathrm{Gap}_{l+1}$ in stage~14. That is a nested chain as in (L2).
\item The gate gas is the restriction of the field's state to it.
\end{itemize}
\begin{proposition}[the fold; proved]\label{prop:fold}
$\relu=\frac12(\mathrm{id}+|\cdot|)$.
\begin{enumerate}[leftmargin=*]
\item The identity half is linear. It commutes with the $O(n)$ symmetry of the weight ensemble, sends quasi-free laws to quasi-free laws,
and creates no cumulants (Proposition~\ref{prop:layeraction}).
\item The fold $|\cdot|$ is the quotient by the sign group $(\Z/2)^n$. Together with the relabelling symmetry $S_n$ of hidden units, this is
the hyperoctahedral group $B_n$, the Weyl group of type~B.
\begin{itemize}
\item In the layer's own frame $z=Wx$ the walls are the reflection hyperplanes of $(\Z/2)^n$, and the cells are its chambers.
\item So the network alternates $O(n)$-equivariant linear maps with $B_n$-equivariant reflecting walls.
\item All births come from the folds.
\end{itemize}
\item At a fold, the even and odd Hermite sectors are Laguerre sectors with half-integer parameters:
\begin{itemize}
\item $\He_{2k}(x)=(-2)^kk!\,L_k^{(-1/2)}(x^2/2)$ and $\He_{2k+1}(x)=(-2)^kk!\,x\,L_k^{(1/2)}(x^2/2)$.
\item The parameters $a=\pm\frac12$, which index Zhou's two level parities through the Jacobi weights $(1-x)^a(1+x)^{-1/2}$, are the
spectral signature of a reflecting wall. Our fold carries the same pair in its Laguerre limit.
\end{itemize}
\item The folded law is an image sum, $p_{|z|}(y)=p(y-\mu)+p(y+\mu)$ for $y>0$. This is the method of images of the reflecting walk.
\end{enumerate}
\end{proposition}

\subsection{The gate gas: arcsine pairs and stars}
\begin{proposition}[proved to leading order; checked]\label{prop:stars}
Let $z\sim N(0,R)$ have unit variances and correlations $\rho_{ij}$, and let $\varepsilon_i=\operatorname{sign}z_i$. Then
$\E\varepsilon_i\varepsilon_j=\frac2\pi\arcsin\rho_{ij}$, and the fourth joint cumulant is
\[
\kappa(\varepsilon_1,\varepsilon_2,\varepsilon_3,\varepsilon_4)=-\frac4{\pi^2}\sum_{i=1}^4\prod_{j\ne i}\rho_{ij}+O(\rho^4).
\]
\end{proposition}
\begin{proof}
Write $\operatorname{sign}=\sum_{k\ \rm odd}a_k\He_k/k!$ with $a_k=2\He_{k-1}(0)\varphi(0)$, so $a_1=\sqrt{2/\pi}$ and $a_3=-\sqrt{2/\pi}$. In the
Gaussian diagram expansion the two-vertex components reproduce the arcsine pairings. The lowest connected diagram with odd degrees on
four vertices is the star of degrees $(3,1,1,1)$, of weight $a_3a_1^3=-(2/\pi)^2$.
\end{proof}
\emph{Check} (orthant probabilities by Genz quadrature): the ratio of the exact cumulant to the star formula is $1.013$, $1.038$, $1.074$
at correlation scales $0.1$, $0.2$, $0.3$. Two readings follow.
\begin{itemize}[leftmargin=*]
\item \emph{Rounding.} The arcsine law is Grothendieck--Krivine rounding of the linear, quasi-free correlation. Tsirelson's theorem
realizes the linear correlations themselves with Clifford observables. The classical gates pay the rounding, up to the Grothendieck
constant in correlation strength.
\item \emph{Stars are the non-quasi-free content.} What the gate gas has beyond pairings begins with stars, gates with three legs. These
are the hub-rooted stars of stage~9 (P2) and the births of the chain. Stage~12 called the memory the ``non-quasi-free remainder''; this
is its first diagram.
\end{itemize}

\subsection{Two sectors, and where the missing information lives}
Every statistic an estimator forms from the instance splits into two parts.
\begin{itemize}[leftmargin=*]
\item An $O(n)$-invariant part: spectral and Gram data, traces, kernels. This is the isotropic sector, and it contains the modular one.
\item A $B_n$-specific remainder, which depends on the neuron basis: Hadamard and diagonal readouts, hub pairings.
\end{itemize}
For the invariant sector, (L1)--(L2) apply in their usual form. Linear spectral statistics of nested rank-one (hub) chains have Gaussian
fluctuations whose covariance is computable from the ensemble, as for Wishart minors. For the $B_n$ sector no such reduction exists.
Corollary~\ref{cor:nofreelunch} says that this is where the information an estimator lacks lives. Lesson (L3) appears as the split
$\relu=\frac12(\mathrm{id}+|\cdot|)$. The identity half carries the first jet, so it transports errors; the fold half creates them.

\section{Complexity: a random circuit with a Gaussian backbone}
\label{sec:complexity}
\subsection{Gaussian backbone, magic folds}
A Gaussian law under linear maps is simulated exactly by $(\mu,C)$. This is the classical counterpart of Gottesman--Knill
(stabilizer states under Clifford gates) and of fermionic linear optics (matchgates). Folds are the only non-Gaussian gates. Their
``magic'' at field $r$ is the energy above the first chaos, which is stage~14's kink law summed over orders:
\[
m(r)=\Var\relu(r+\xi)-\Phi(r)^2=\sum_{k\ge2}\frac{\He_{k-2}(r)^2\varphi(r)^2}{k!},\qquad m(0)=\tfrac12-\tfrac1{2\pi}-\tfrac14=0.0908 .
\]
\begin{proposition}[the magic budget; checked]
$M_\infty=\int m(r)\,dr=0.18806$. The order $k$ contributes $\int E_k=0.141,0.024,0.009,0.004,\dots$ for $k=2,3,4,5$. Under stage~12's field law
$r\sim N(0,t_l)$, the magic per neuron is $\E\,m(r)\simeq M_\infty/\sqrt{2\pi(1+t_l)}$:
\begin{itemize}[leftmargin=*]
\item $0.0577$ against $0.0531$ at $t=1$;
\item $0.0346$ against $0.0336$ at $t=4$;
\item $0.0211$ against $0.0208$ at $t=12$.
\end{itemize}
So a late layer of the competition networks injects about $0.021\,n\approx22$ units, and the total over depth grows like $\log\tau$
(stage~12's harmonic Fatou law).
\end{proposition}
For circuits dominated by Clifford gates, the cost of exact simulation is exponential in the magic. The mean, however, needs only
connected diagrams of low order. The cumulant (birth) expansion is a low-magic expansion. Its cost is set by how many magic insertions
must be tracked \emph{coherently}, which for a chain is its number of live $\kappa_3$ families. Theorem~\ref{thm:cocycle} says that
incoherent magic costs only its energy.

\subsection{The Pauli-propagation dictionary}
\begin{center}\small
\begin{tabular}{>{\raggedright\arraybackslash}p{5.6cm}>{\raggedright\arraybackslash}p{8.4cm}}\toprule
random quantum circuit & bias-free He \textsc{ReLU} network\\\midrule
Pauli weight of a path & jet order $k$ (Hermite degree at a wall)\\
random single-qubit gate & a fresh Gaussian row (Theorem~\ref{thm:rowmehler})\\
paths orthogonal on average & births orthogonal across rows and layers (Theorem~\ref{thm:cocycle})\\
noise damps weight $k$ by $(1-p)^k$ & jets $k\ge2$ live at walls: $\E_r[s^{2k-2}\J_k^2]\propto(1+t)^{-1/2}$, algebraic, not exponential\\
weight $1$ propagates & the first jet $\Phi(r)$ transports with factor $0.65$--$0.94$, near-critical\\
truncation MSE $=$ weight of dropped paths & the cocycle sum rule, measured to $4$--$10\%$\\
the one non-orthogonal coupling & the dilation (radial norm)\\\bottomrule
\end{tabular}
\end{center}
The network is a weakly damped, near-critical random circuit. Damping is algebraic, through the wall density, so memory persists over
many layers (stage~9: eleven ages matter). Births are not damped where they are born, so the late layers dominate.

\subsection{Calibration against quantum mean estimation}
Take $\sigma^2=0.0787$ (average output variance), one forward pass $=2Ln^2$, and $B=2048n^3$. Quantum mean estimation needs about
$\sigma/\varepsilon$ coherent passes per output, and about $\sqrt n\,\sigma/\varepsilon$ for all $n$ outputs. The latter is the multi-observable
bound. We ignore reversibility and error-correction overheads, which is optimistic for the quantum side.
\begin{center}\footnotesize\setlength{\tabcolsep}{4pt}
\begin{tabular}{lcccc}\toprule
target raw MSE & classical Monte Carlo & quantum, one output & quantum, all outputs & deterministic closures\\\midrule
$1.56\times10^{-8}$ & $77\,B$ & $0.034\,B$ & $1.10\,B$ & about $0.2\,B$ (v56)\\
$1.0\times10^{-9}$ & $1201\,B$ & $0.135\,B$ & $4.33\,B$ & ---\\\bottomrule
\end{tabular}
\end{center}
For all outputs, the closures beat generic quantum mean estimation by about $5\times$. They exploit structure that a generic algorithm does
not: one covariance gives every output at once, the Gaussian backbone. Large further gains can only come from more structure. They
cannot come from better generic sampling, classical or quantum.

\subsection{The low-degree view in weight space}
Average the MSE over networks. In $L^2$ of the weight ensemble the truth and any estimator have Hermite (weight-chaos) expansions, and
errors of different weight-degree are orthogonal. Corollary~\ref{cor:nofreelunch} says the following in this language.
\begin{itemize}[leftmargin=*]
\item Corrections built from features of degree $0$ in the anisotropic directions of a fresh row reach only the isotropic births.
\item An anisotropic birth at row $i$ is a polynomial of degree $k\ge1$ in those directions, contracted with the instance's
$\delta\kappa_k$.
\item The hardness of the problem is the number of anisotropic cumulant directions that must be resolved at the late layers. Global
calibration cannot change it.
\end{itemize}

\subsection{Coding}
\begin{itemize}[leftmargin=*]
\item \emph{A layer is a one-bit embedding.} $E[d_H(\operatorname{sign}Wx,\operatorname{sign}Wy)]=n\,\theta(x,y)/\pi$. Each layer one-bit-senses the
previous one. The closure's two-point law is the code's distance law, which is the stage-12 angle map.
\item \emph{Two Grams of one operator.} Let $(\cG f)_i=\E[\bar g_i f]$ with centred gates $\bar g_i$.
\begin{itemize}
\item $\cG\cG^*$ is the neuron-side gate covariance, which the closure computes.
\item $\cG^*\cG$ is the input-side code kernel $\sum_i\bar g_i(x)\bar g_i(y)$, which NNGP theory computes.
\item They have the same non-zero spectrum. The neuron picture and the input picture are dual codes of one operator, a
MacWilliams-type duality.
\end{itemize}
\item \emph{Lattice side.} The cell masses are the code's weight distribution (stage~14's gap labels). The tight frames of contact
normals are its local codes.
\end{itemize}

\section{What a theory-native estimator computes}
\label{sec:system}
We now derive the estimator from Sections~\ref{sec:hopf}--\ref{sec:complexity}, and from nothing else.

\subsection{Three sectors with three different statuses}
By Theorem~\ref{thm:cocycle}, an estimator's error has three parts, and each needs a different treatment.
\begin{description}[leftmargin=*]
\item[S1, the modular sector: compute exactly.] The dilation is the one coherent mode. The first two jets carry it with factor exactly
$1$, and its errors add in amplitude across layers. It was $61\%$ of the Gaussian closure's MSE. It must not be approximated or fitted. Use
stage~13's dilation package (exact to $O(\varepsilon^2)$ on scale mixtures), or stage~14's radial partition function, at $O(n^2)$ per
layer.
\item[S2, the linear cocycle: apply exactly.] $T_l=\Phi(r_l)\odot W_l$ is a matrix--vector product. Its adjoint gives each row's downstream
gain $a_{l,i}$, the norm of row $i$ of $T_L\cdots T_{l+1}$.
\item[S3, the births: the only thing approximated.]
\[\beta_{l,i}=\sum_k\J_k(r_{l,i})\langle\delta\kappa_k,w_{l,i}^{\otimes k}\rangle/k!\,.\]
Compute them
with true-law jets \emph{and} the matching hyperedges, at one consistent order (Corollary~\ref{cor:consistency}), and only where
birth $\times$ gain is large.
\end{description}
This is the interaction picture. The Gaussian backbone is the free field. Non-Gaussianity enters only as births, and their effect on
later means is linear and exact.

\subsection{The instance decides the estimator}
The question was how an optimal system could fail to couple the sampled weights to the shape of its own computation. In the cocycle
picture the coupling is forced, and it can be computed.
\begin{proposition}[additivity makes allocation a knapsack]\label{prop:knapsack}
Under the orthogonality of Theorem~\ref{thm:cocycle}, the expected MSE of an estimator that captures the fraction $c_{l,i}$ of the birth at
$(l,i)$ is
\[
\sum_{l,i}(1-c_{l,i})^2\,a_{l,i}^2\,\E\beta_{l,i}^2 ,\qquad \E\beta_{l,i}^2\approx\sum_{k\ge2}\frac{(s^{k-1}\J_k(r_{l,i}))^2}{k!}\,\epsilon_{l,k}^2 ,
\]
where $\epsilon_{l,k}$ is the relative row-projected error of the order-$k$ cumulant at layer $l$. The terms are additive. So the optimal
allocation of a FLOP budget over (layer, row, order) is a fractional knapsack, solved greedily by captured energy per FLOP.
\end{proposition}
The weights come from a cheap first pass that depends on the instance: the fields $r_{l,i}$ and scales $s_{l,i}$ from the Gaussian
backbone, and the gains $a_{l,i}$ from one adjoint sweep. Both are $O(n^3)$ once. The second pass spends the budget where the weights
say. Different networks then get different estimators, as they should: a network with more near-wall rows at depth gets more late-layer
effort. Dynamically, the allocation of layer $l+1$ can be revised once layer $l$'s births are known.

\subsection{Where the budget goes on the official networks}
On network~0 the exact first-order chain's births are $5\%$ in layers $1$--$6$ and $81\%$ in $11$--$16$.
\begin{itemize}[leftmargin=*]
\item \textbf{Last-layer channels.} A fit of the last layer's residual-variance profile by squared jets, weighted over the actual
fields, attributes $59\%$ of the last-layer birth to the second jet (a covariance error of $h_{15}$, profile $\varphi^2$). The third jet
($\kappa_3$, profile $r^2\varphi^2$) takes $31\%$ and the fourth $11\%$. The fit is coarse: $20$ bins, residual $0.30$.
\item \textbf{What causes them.} By Theorem~\ref{thm:walljet}, covariance births at a late layer come from the $(2,2)$ and $(3,1)$
hyperedges and the $\kappa_3$-pair diagrams in $\Cov(\relu,\relu)$. These are what the exact chain lacks and what v56's $\kappa_4$ sector
partly adds. The theory points the effort at late covariance births.
\end{itemize}
\paragraph{Heisenberg reads for the late layers.} The late readouts dominate, so read the old $\kappa_3$ families \emph{backwards}.
\begin{itemize}[leftmargin=*]
\item \emph{The read.} Transport the readout rows of a late layer $L'$ to the layer $l_f$ where a family was born, giving
$V=(B_{L'-1}\cdots B_{l_f+1})^\top w$. Then evaluate $\sum_ac_a(V u_a)(Vv_a)^2$ exactly. Old sources need no compression, and stage~9 showed
that Frobenius compression fails against Hadamard readouts.
\item \emph{Restricting rows.} Only near-wall rows need this ($\J_3\propto r\varphi(r)$): about half at $t\approx12$.
\item \emph{Cost.} About $6$ late readout layers, each with backward transport and family reads over at most $16$ layers on half the
rows, cost on the order of $100\,n^3\approx0.05\,B$.
\item \emph{Early layers} keep a short forward window. Their births are small and discounted by the gain.
\end{itemize}
This is a meet-in-the-middle estimator: Schrödinger (forward) where births are cheap and discounted, Heisenberg (adjoint) where they
dominate.

\subsection{What this system is not}
\begin{center}\small
\begin{tabular}{>{\raggedright\arraybackslash}p{6.4cm}>{\raggedright\arraybackslash}p{7.6cm}}\toprule
standard (public chain, v56) & theory-native\\\midrule
one closure, the same at every layer and for every network & three sectors; allocation decided per instance (Proposition~\ref{prop:knapsack})\\
law propagator with feedback of corrections into $(\mu,C)$ & interaction picture: exact backbone, births transported by the exact cocycle\\
dilation absorbed in fitted pair classes and counterterms & modular sector computed exactly; counterterms proved useless beyond the isotropic part\\
old sources compressed forward & late readouts done backward (Heisenberg), exactly, on near-wall rows\\
Gaussian jets with Edgeworth patches & true-law jets and hyperedges at one consistent order\\\bottomrule
\end{tabular}
\end{center}

\subsection{Predictions and the next measurements}
\begin{enumerate}[leftmargin=*]
\item \textbf{The modular sector alone.} Making it exact in the Gaussian closure removes at most its collective share, $61\%$ of the MSE. The
MSE then goes from $4.1\times10^{-6}$ towards $1.6\times10^{-6}$, at $O(n^2)$ cost.
\item \textbf{Counterterms.} On top of an exact first-order chain, counterterms of any kind gain at most about
$0.21\times0.07+0.01\approx2.5\%$.
\item \textbf{The last layer.} The true covariance of $h_{15}$ will show that about $60\%$ of the exact chain's last-layer birth is
row-projected covariance error. Testing this needs one Monte Carlo run of $\Cov(h_{15})$, about $2\times10^6$ samples.
\item \textbf{Early precision.} Readout and birth precision in layers $1$--$6$ can be relaxed about $3\times$ in amplitude at a cost of $\lesssim5\%$
of the MSE, provided the sources born there are kept for the late Heisenberg reads.
\item \textbf{Biased networks} (stage~14's conjecture). With biases the dilation is no longer free. The collective share and the
coherent part of the cocycle should then shrink with the bias scale.
\end{enumerate}

\section{Checks}
\label{sec:verify}
The script \texttt{scripts/\allowbreak verify\_\allowbreak stage15.py} runs in under a minute. Its transcript is
\texttt{notes/\allowbreak stage15/\allowbreak verify\_\allowbreak stage15.txt}. Section~2 needs
the official network~0 and its $10^9$-sample truth. The exact first-order chain is the organisers' $K=3$ simple chain with radial
$\kappa_4$, ported and verified in stage~9 (MSE $3.53\times10^{-8}$).
\begin{center}\small
\begin{tabular}{>{\raggedright\arraybackslash}p{5.2cm}>{\raggedright\arraybackslash}p{8.8cm}}\toprule
statement & result\\\midrule
connected-cumulant law $\E[\Ap_p\Ap_q]$, $p,q\le7$ & max relative difference $2.3\times10^{-13}$\\
Appell expansion of $\Cov(X^3-2X,Y^4+Y)$ & $2.3159134053$ both ways\\
wall-jet expansion of $\Cov(\relu,\relu)$ & errors $-4.7\times10^{-2}\to+4.0\times10^{-4}$ (orders $1\to7$), Gaussian closure $-1.2\times10^{-2}$\\
consistency (jets without hyperedges) & $-2.7\times10^{-2}$, worse than the Gaussian closure\\
row Mehler, $p,q\le3$ & $|z|\le1.4$ over $4\times10^5$ rows; radial leak $0.2346$ vs $0.2364$\\
first-jet share of the final MSE & $0.953$ (Gaussian closure), $0.790$ (exact chain)\\
cocycle sum rule & $2.59/4.06$ (Gaussian), $1.50/1.57$ after removing the collective mode, $3.20/3.53$ (exact chain)\\
collective share of inherited error & $0.48$--$0.68$ (Gaussian), $0.00$--$0.03$ (exact chain)\\
births in layers $11$--$16$ & $51\%$ (Gaussian), $81\%$ (exact chain)\\
last-layer channels (exact chain) & covariance $0.59$, $\kappa_3$ $0.31$, $\kappa_4$ $0.11$; field-explained part of the remainder $0.073$\\
gate-gas star cumulant & exact/star $=1.013$, $1.038$, $1.074$ at scales $0.1$, $0.2$, $0.3$\\
magic $M_\infty$ and the layer law & $0.18806$; $\E m$ vs $M_\infty/\sqrt{2\pi(1+t)}$: $0.0577/0.0531$, $0.0346/0.0336$, $0.0211/0.0208$\\
quantum calibration & Monte Carlo $77\,B$; quantum $0.034\,B$ (one output), $1.10\,B$ (all outputs) at MSE $1.56\times10^{-8}$\\\bottomrule
\end{tabular}
\end{center}
What is not checked here:
\begin{itemize}[leftmargin=*]
\item the Stein cross-term bound in Theorem~\ref{thm:cocycle}, beyond the measured sum rules;
\item the knapsack allocation and the Heisenberg reads, which are designs, not results;
\item the Laguerre parallel with Zhou's parities, which is structural, not quantitative.
\end{itemize}

\end{document}
